\documentclass[11pt,a4paper]{article}
\usepackage{gtu-el-yazisi}
\begin{document}
\notbilgi{Görüntü İşleme}{Görüntü eşikleme (thresholding)}{Ders8.pdf}
\textit{Not: Orijinal dosyadaki basılı sayfalar (s. 2) dijitalleştirilmedi; yalnızca el yazısı sayfalar aktarıldı.}

\sayfa{1}
\section{Otomatik eşikleme}
\sekilnotu[1]{Basılı kitap şekli (Şekil 10.38: gürültülü parmak izi, histogram, global eşik ile segmentasyon sonucu); dijitalleştirilmedi.}

Tüm görüntüye aynı eşikse global eşikleme

bölgesel farklı eşikler local eşikleme

\subsection{Global eşikleme}
\[
g(x,y)=\begin{cases} 1 & \text{if } f(x,y)>T\\ 0 & \text{o.w} \end{cases}
\]

\section{Basic Global Tresholding}
\begin{enumerate}[label=\Alph*), leftmargin=*]
\item Initialize $T_0=\dfrac{1}{MN}\displaystyle\sum_x\sum_y f(x,y)$, $i=0,\ \epsilon$
\item $i\leftarrow i+1$
\item Segment the image
  \begin{align*}
  (x,y)&\in G_1 \quad\text{if } f(x,y)\le T_{i-1}\\
  (x,y)&\in G_2 \quad\text{if } f(x,y)> T_{i-1}
  \end{align*}
\item Compute $\mu_1,\mu_2$
  \begin{align*}
  \mu_1&=\frac{1}{\underbrace{k_1}_{\text{piksel sayısı}}}\sum_x\sum_y f(x,y), \qquad x,y\in G_1\\
  \mu_2&=\frac{1}{k_2}\sum_x\sum_y f(x,y), \qquad x,y\in G_2\\
  k_1+k_2&=M\cdot N
  \end{align*}
\item compute new Treshold \quad $T_i=\dfrac{1}{2}(\mu_1+\mu_2)$
\item if $|T_i-T_{i-1}|>\epsilon$ \quad go to B)
\item $T=T_i$
\end{enumerate}

\sayfa{3}
\sekilnotu[3]{Basılı kitap şekilleri (Şekil 10.47 ve 10.49: alt görüntülerin histogramları; Otsu ve hareketli ortalama ile eşikleme sonuçları); dijitalleştirilmedi.}

\section{Variable Tresholding based on local image processing}
\elcizimi[0.4]{sekiller/Ders8-s3-1.png}[3]

\textcolor{kirmizikalem}{Set Treshold} \quad $T_{x,y}=a\,\sigma_{x,y}+M_{x,y}$, \quad $a>0$

\elcizimi[0.45]{sekiller/Ders8-s3-2.png}[3]

Impulse noise tuzbiber gürültüsü

\sayfa{4}
\sekilnotu[4]{Basılı kitap şekilleri (Şekil 10.50: sinüzoidal gölgelenmiş metin görüntüsü, Otsu ve hareketli ortalama ile eşikleme; Şekil 10.41: gürültülü görüntü, histogramı ve Otsu sonucu); dijitalleştirilmedi.}

\okunamadi{Alıray} çekiyoruz

\elcizimi[0.7]{sekiller/Ders8-s4-1.png}[4]

\sayfa{5}
Histogram hesabı yaparken sadece edge'den alacağız

% (şekil içindeki not: sadece inner ve outer edge'in piksellerini)
\elcizimi[0.75]{sekiller/Ders8-s5-1.png}[5]

\sekilnotu[5]{Basılı kitap şekli (Şekil 10.42: gradyan büyüklüğü eşiklenerek elde edilen kenar piksellerinin histogramı ve Otsu eşikleme sonucu); dijitalleştirilmedi.}

Check görüntüsü var 2. baskıdan

\sayfa{6}
\section{Otsu Tresholding}
\elcizimi[0.55]{sekiller/Ders8-s6-1.png}[6]
\elcizimi[0.6]{sekiller/Ders8-s6-2.png}[6]

\section{Normalized histogram $p(i)$}
$i=1,\dots,I$\\
$I$ is number of graylevels
\[ \sum_{i=1}^{I} p(i)=1 \]
\[ \mu_k(t) \qquad k=1,2 \]
\[ \sigma_k^2(t) \]
\begin{align*}
\mu_1(t)&=\sum_{i=1}^{t} i\,\frac{p(i)}{q_1(t)} & q_1(t)&=\sum_{i=1}^{t}p(i)\\
\mu_2(t)&=\sum_{i=t+1}^{I} i\cdot\frac{p(i)}{q_2(t)} & q_2(t)&=\sum_{i=t+1}^{I}p(i)\\
& & q_1(t)+q_2(t)&=1
\end{align*}
$\leftarrow$ $p$'lerin yüzde kaçı $t$'nin altında
\begin{align*}
\sigma_1^2(t)&=\sum_{i=1}^{t}\bigl(i-\mu_1(t)\bigr)^2\,\frac{p(i)}{q_1(t)}\\
\sigma_2^2(t)&=\sum_{i=t+1}^{I}\bigl(i-\mu_2(t)\bigr)^2\,\frac{p(i)}{q_2(t)}
\end{align*}

\section{within group variance}
\[ \sigma_w^2(t)=q_1(t)\,\sigma_1^2(t)+q_2(t)\,\sigma_2^2(t) \]
$\nearrow$ ortalama variance hesabı

Otsu : minimize $\sigma_w^2(t)$

\sayfa{7}
\elcizimi[0.5]{sekiller/Ders8-s7-1.png}[7]

görüntünün Varyansı
\begin{align*}
\sigma^2&=\sum_{i=1}^{I}(i-\mu)^2\,p(i)\\
\mu&=\sum_{i=1}^{I} i\,p(i)
\end{align*}
\[
\underbrace{\sigma^2}_{\text{constant}}=\sigma_w^2(t)+\underbrace{q_1(t)\bigl[\mu_1(t)-\mu\bigr]^2+q_2(t)\bigl[\mu_2(t)-\mu\bigr]^2}_{\substack{\sigma_B^2(t)\ \ \text{maximize}\\ \rightarrow\ \text{between group variance}}}
\]

\section{Otsu:}
\begin{align*}
\mu&=\sum i\,p(i)\\
q_1(0)&=0\\
\mu_1(0)&=0
\end{align*}
\[ \text{find }\arg\max_{t=1,\dots,I}\ \sigma_B^2(t) \]
\begin{align*}
q_1(t)&=q_1(t-1)+p(t)\\
q_2(t)&=1-q_1(t)\\
\mu_1(t)&=\bigl[q_1(t-1)\,\mu_1(t-1)+t\,p(t)\bigr]/q_1(t)\\
\mu_2(t)&=\bigl[\mu-q_1(t)\,\mu_1(t)\bigr]/q_2(t)\\
\sigma_B^2(t)&=q_1(t)\,q_2(t)\,\bigl[\mu_1(t)-\mu_2(t)\bigr]^2
\end{align*}

\section{Canny edge detection}
\begin{enumerate}[label=\arabic*., leftmargin=*]
\item Smooth the image:
\begin{align*}
g(x,y)&=G_\sigma(x,y)\ast f(x,y)\\
G_\sigma(x,y)&=\frac{1}{\sqrt{2\pi\sigma^2}}\,e^{-\frac{x^2+y^2}{\sigma^2}}
\end{align*}
\end{enumerate}

\sayfa{8}
\begin{enumerate}[label=\arabic*., start=2, leftmargin=*]
\item $\displaystyle M(x,y)=\sqrt{g_x^2(x,y)+g_y^2(x,y)}$
\[ \alpha(x,y)=\tan^{-1}\bigl(g_y(x,y)/g_x(x,y)\bigr) \]
\item Treshold $M$:
\[ M_T(x,y)=\begin{cases} M(x,y) & \text{if } M(x,y)>T\\ 0 & \text{o.w.}\end{cases} \]
\item Suppress non-maxima pixels in $M_T$ \quad (kalın olmaması için edgelerin)

\elcizimi[0.5]{sekiller/Ders8-s8-1.png}[8]

\[ 0\ 0\ 0 \quad 150 \quad 160 \quad \underbrace{170}_{\text{bunu al}} \quad 162 \quad 142 \quad 0\ 0\ 0\ 0 \]
\item İki eşik bakıyoruz

$z_1 \qquad z_2 \qquad\qquad z_1<z_2$

Treshold $M_T$ by $z_1$ and $z_2$ to find $T_1, T_2$

Canny'nin hedefi kırık kenar olmayacak.

% (şekil içindeki not: buna göre sol bakıyor, varsa devam ediyor)
\elcizimi[0.8]{sekiller/Ders8-s8-2.png}[8]
\end{enumerate}

\end{document}
